\section{Piles électrochimiques et alimentation d'une voiture télécommandée}
\textbf{Source du sujet~:} Baccalauréat STL Physique-Chimie, Polynésie – 2023.

\subsection{Introduction}
Une pile Daniell est un dispositif électrochimique réalisant la conversion d'énergie chimique
en énergie électrique par le biais d'une réaction d'oxydoréduction. Le principe de
fonctionnement de cette pile a été démontré par Edmond Becquerel et a été perfectionné
par le chimiste britannique John Daniell en 1836, au moment où le développement du
télégraphe faisait apparaître un besoin urgent de sources de tension sûres et constantes.

La pile Daniell vient corriger certains défauts de la pile Volta : elle est simple de
construction, commode d'usage et de tension stable durant son utilisation, si bien qu'elle
servit pendant longtemps de pile étalon en laboratoire.

On appelle pile alcaline un type de pile électrique dont l'électrolyte est alcalin, c’est-à-dire basique. L’un des modèles les plus courants est la pile alcaline zinc-dioxyde de manganèse (Zn-MnO2) que l’on utilise pour alimenter des jouets par exemple (\emph{Source
Wikipédia}).

L’objectif de cet exercice est d’étudier une pile Daniell réalisée au laboratoire puis d’étudier les performances d’une voiture télécommandée, alimentée à l’aide de piles alcalines.

\subsection{Données}
\begin{itemize}
    \item Masses molaires atomiques : $M(\mathrm{Cu}) = \SI{63.5}{\gram\per\mole}$ ; $M(\mathrm{Zn}) = \SI{65.4}{\gram\per\mole}$ ;
    \item Charge élémentaire : $e = 1,60 \times 10^{-19}\ \mathrm{C}$ ;
    \item Constante d'Avogadro : $N_{\mathrm{A}} = 6,02 \times 10^{23}\ \mathrm{mol}^{-1}$ ;
    \item Couples oxydant/réducteur de la pile Daniell : $\mathrm{Zn^{2+}(aq) / Zn(s)}$ et $\mathrm{Cu^{2+}(aq) / Cu(s)}$ ;
    \item Couples oxydant/réducteur de la pile alcaline : $\mathrm{ZnO(s) / Zn(s)}$ et $\mathrm{MnO_2(s) / MnO(OH)(s)}$ ;
    \item Capacité d'une pile alcaline : de 2 000 mAh à 3 000 mAh, soit respectivement de 7 200 C à 10 800 C ;
    \item Masse de la voiture radiocommandée : $m = \SI{741}{\gram}$.
\end{itemize}

\subsection{Étude d'une pile Daniell}
Une pile Daniell est constituée d'une lame de zinc plongée dans une solution contenant un volume $V_{0} = \SI{20}{\milli\liter}$ d'une solution de sulfate de zinc ($\mathrm{Zn^{2+}(aq) + SO_4^{2-}(aq)}$) de concentration $C_0 = \SI{1.0e-2}{\mole\per\liter}$ et d'une lame de cuivre plongée dans une solution contenant un volume $V_{0} = \SI{20}{\milli\liter}$ de solution de sulfate de cuivre ($\mathrm{Cu^{2+}(aq) + SO_4^{2-}(aq)}$) de concentration $C_0 = \SI{1.0e-2}{\mole\per\liter}$. Les deux solutions sont reliées par un pont salin contenant une solution de nitrate d'ammonium ($\mathrm{NH_4^+(aq) + NO_3^-(aq)}$).

Un voltmètre est branché aux bornes de la pile pour mesurer la valeur de la tension à vide $U_{0}$ lorsque la pile ne débite pas de courant électrique. Le voltmètre mesure une tension à vide positive~: $U_{0} = \SI{1.102}{\volt}$.

\begin{questionmult}{pileD1}
    Préciser, en justifiant la réponse, la polarité de chaque électrode.
    \begin{EnvQuadrillage}[NbCarreaux=10x4,Grille=Seyes,Marge=1]
    \end{EnvQuadrillage}
    \AMCOpen{}{
        \mauvaise[F]{NA}\scoring{b=0}
        \bonne[A1]{Polarité correcte}\scoring{b=1}
        \bonne[A2]{Justification correcte}\scoring{b=1}
    }
    \explain{
    La tension mesurée est positive, ce qui indique que l'électrode de cuivre est la borne positive et l'électrode de zinc la borne négative.
    }
\end{questionmult}

\begin{questionmult}{pileD2}
    Écrire les équations de demi-réaction électroniques se produisant à l'anode et à la cathode. En déduire la réaction d'oxydoréduction modélisant la transformation chimique au sein de la pile.
    \begin{EnvQuadrillage}[NbCarreaux=10x4,Grille=Seyes,Marge=1]
    \end{EnvQuadrillage}
    \AMCOpen{}{
        \mauvaise[F]{NA}\scoring{b=0}
        \bonne[A1]{Équations correctes}\scoring{b=1}
        \bonne[A2]{Réaction correcte}\scoring{b=1}
    }
    \explain{
    Les demi-équations sont :
    $$
    \ce{Zn(s) -> Zn^{2+}(aq) + 2 e^{-}}
    $$
    $$
    \ce{Cu^{2+}(aq) + 2 e^{-} -> Cu(s)}
    $$
    La réaction globale est :
    $$
    \ce{Zn(s) + Cu^{2+}(aq) -> Zn^{2+}(aq) + Cu(s)}
    $$
    }
\end{questionmult}

\begin{questionmult}{pileD3}
    Donner le sens de déplacement des ions ammonium $\mathrm{NH_4^+}$ (aq) et des ions nitrate $\mathrm{NO_3^-}$ (aq) dans le pont salin.
    \begin{EnvQuadrillage}[NbCarreaux=10x4,Grille=Seyes,Marge=1]
    \end{EnvQuadrillage}
    \AMCOpen{}{
        \mauvaise[F]{NA}\scoring{b=0}
        \bonne[A1]{Sens correct}\scoring{b=1}
    }
    \explain{
    Les ions ammonium $\mathrm{NH_4^+}$ se déplacent vers la solution de sulfate de cuivre, et les ions nitrate $\mathrm{NO_3^-}$ se déplacent vers la solution de sulfate de zinc.
    }
\end{questionmult}

\begin{questionmult}{pileD4}
    Déterminer le réactif limitant de la réaction.
    \begin{EnvQuadrillage}[NbCarreaux=10x4,Grille=Seyes,Marge=1]
    \end{EnvQuadrillage}
    \AMCOpen{}{
        \mauvaise[F]{NA}\scoring{b=0}
        \bonne[A1]{Réactif correct}\scoring{b=1}
        \bonne[A2]{Calcul correct}\scoring{b=1}
    }
    \explain{
    La quantité de matière initiale de zinc est $n_{\mathrm{Zn}} = \frac{m_{\mathrm{Zn}}}{M(\mathrm{Zn})} = \frac{50.2}{65.4} = \SI{0.767}{\mole}$.
    La quantité de matière initiale de cuivre est $n_{\mathrm{Cu}} = \frac{m_{\mathrm{Cu}}}{M(\mathrm{Cu})} = \frac{62.8}{63.5} = \SI{0.989}{\mole}$.
    Le réactif limitant est le zinc.
    }
\end{questionmult}

\begin{questionmult}{pileD5}
    En déduire la capacité $Q_D$ de la pile Daniell. Commenter la valeur obtenue.
    \begin{EnvQuadrillage}[NbCarreaux=10x4,Grille=Seyes,Marge=1]
    \end{EnvQuadrillage}
    \AMCOpen{}{
        \mauvaise[F]{NA}\scoring{b=0}
        \bonne[A1]{Calcul correct}\scoring{b=1}
        \bonne[A2]{Valeur correcte}\scoring{b=1}
    }
    \explain{
    La quantité d'électricité est donnée par $Q_D = n_{\mathrm{e^{-}}} \times F = 2 \times n_{\mathrm{Zn}} \times F = 2 \times 0.767 \times 96500 = \SI{1.48e5}{\coulomb}$.
    }
\end{questionmult}

\subsection{Étude d'une pile alcaline Power+}
Les piles alcalines sont fabriquées sous forme de cylindres et de boutons. Une pile cylindrique est contenue dans un tube d'acier, qui sert de collecteur du courant de la cathode.

\begin{questionmult}{pileA1}
    Identifier en justifiant la réponse, le réactif limitant de la pile alcaline.
    \begin{EnvQuadrillage}[NbCarreaux=10x4,Grille=Seyes,Marge=1]
    \end{EnvQuadrillage}
    \AMCOpen{}{
        \mauvaise[F]{NA}\scoring{b=0}
        \bonne[A1]{Réactif correct}\scoring{b=1}
        \bonne[A2]{Justification correcte}\scoring{b=1}
    }
    \explain{
    Le réactif limitant est le zinc, car il est entièrement consommé avant le dioxyde de manganèse.
    }
\end{questionmult}

\subsection{Alimentation d'une voiture radiocommandée et étude des performances}
Une voiture télécommandée est alimentée par un bloc de 6 piles alcalines AA Power+ étudiées précédemment.

\begin{questionmult}{voiture1}
    Déterminer la valeur de la durée d'utilisation maximale attendue de la voiture télécommandée. Commenter le résultat.
    \begin{EnvQuadrillage}[NbCarreaux=10x4,Grille=Seyes,Marge=1]
    \end{EnvQuadrillage}
    \AMCOpen{}{
        \mauvaise[F]{NA}\scoring{b=0}
        \bonne[A1]{Calcul correct}\scoring{b=1}
        \bonne[A2]{Valeur correcte}\scoring{b=1}
    }
    \explain{
    La durée d'utilisation maximale est donnée par $t = \frac{Q_B}{I} = \frac{10000}{0.6} = \SI{1.67e4}{\second} = \SI{4.64}{\hour}$.
    }
\end{questionmult}
